\section*{2 Preliminaries}
For a positive integer $n$, we let $[n]$ denote the subset of integers $\{1,\ldots,n\}$.
\subsection*{2.1 DeMorgan Formulas and Formula Size}
Our notion of formulas will be DeMorgan formulas. A DeMorgan formula $\phi$ on $n$-inputs of size $s$ is given by:
\begin{itemize}
\item a directed rooted binary tree on the vertex set $[2s-1]$, specified by a subset $E_\phi\subseteq[2s-1]\times[2s-1]$ of edges, and
\item a gate labeling function $\tau_\phi:[2s-1]\to\{\land,\lor\}\cup\{0,1,x_1,\ldots,x_n,\neg x_1,\ldots,\neg x_n\}$
\end{itemize}
where $\tau$ takes values in $\{\land,\lor\}$ on the internal nodes in $\phi$ and $\tau$ takes values in $\{0,1,x_1,\ldots,x_n,\neg x_1,\ldots,\neg x_n\}$ on the leaf nodes in $\phi$. The edges in $E_\phi$ point from inputs towards outputs. We note that our definition implicitly uses the fact that a binary tree with $s$ leaf nodes has $s-1$ internal nodes. We also note that in our definition we do not need to specify the ``left'' and ``right'' child of an internal node since our gate set $\{\land,\lor\}$ is made up of symmetric functions. We will define a notion of formula isomorphism in Section 2.2.
We will use the notation $|\phi|$ to denote the size of a formula $\phi$ (i.e. the number of leaves in the binary tree underlying $\phi$). Given a Boolean function $f$, we denote the minimum formula size of $f$ by
\[L(f)=\min\{|\phi|:\phi\text{ is a formula computing }f\}.\]
We say a formula $\phi$ is an optimal formula for a Boolean function $f$, if $\phi$ computes $f$ and $|\phi|=L(f)$.
We note, however, that all of our results except the ones presented in Section 7 apply equally well to formulas with arbitrary fan-in-two gates (i.e. the formulas over the $B_2$ basis). Moreover, all our results are hold for other size notions such as gates and wires.
\subsection*{2.2 Optimal Formulas and Formula Isomorphism}
Since our results will depend on the number of formulas satisfying certain properties, we will be clear about when exactly we are saying formulas are distinct in our count.
In particular, as we have defined formulas, one can obtain many optimal formulas from a single optimal formula by relabeling the nodes in underlying binary tree.
Thus, it will be useful to define an isomorphism on formulas and only count formulas modulo this isomorphism. In particular, we will define two formulas to be isomorphic if they are isomorphic as labelled binary trees.
In order to properly define this, we introduce some notation. If $\phi$ is a formula of size $s$ with an underlying edge set $E_\phi$ and a labelling function $\tau_\phi$ and $\sigma:[2s-1]\to[2s-1]$ is a permutation, then we let $\psi=\sigma(\phi)$ be the formula of size $s$ whose edge set $E_\psi$ is given by
\[E_\psi=\{(\sigma(i),\sigma(j)):(i,j)\in E_\phi\}\]
and whose labelling function $\tau_\psi$ is given by $\tau_\psi(\sigma(i))=\tau_\phi(i)$.
We say two formulas $\phi$ and $\phi'$ are isomorphic if $|\phi|=|\phi'|$ and there is a permutation $\sigma$ such that $\phi'=\sigma(\phi)$.
From each equivalence class of isomorphic formulas, we pick a single representative that we call the canonical formula for that equivalence class. Note that for our purposes we do not need that this canonical formula to be computable, as we will just be using them in our analysis. Then we define $\Canon{k}(f)$ to be the set of canonical formulas that are optimal for computing $f$ up to an additive $k$-term. In other words
\[\Canon{k}(f)=\{\phi:\phi\text{ is a canonical formula and }|\phi|\leq L(f)+k\}.\]
\subsection*{2.3 MFSP, Search-MFSP and Conventions on $n$ and $N$}
We now define the Minimum Formula Size Problem denoted MFSP.
\setcounter{nttheorem}{9}
\begin{ntdefinition}[Author Definition 10: MFSP]
We define the problem MFSP as follows
\begin{description}\item[Given:] a truth table of a Boolean function $f$ and an integer size parameter $s\geq 1$\item[Determine:] if $L(f)\leq s$.\end{description}
\end{ntdefinition}
We define the search version of MFSP analogously.
\begin{ntdefinition}[Author Definition 11: Search-MFSP]
Search-MFSP is the problem defined as follows:
\begin{description}\item[Given:] a truth table of a Boolean function $f$\item[Output:] a formula $\phi$ of size $L(f)$ computing $f$.\end{description}
\end{ntdefinition}
We note that MFSP$\in\mathrm{NP}$ since given a minimum-sized formula as a witness, one can check that this indeed computes $f$ efficiently since the truth table of $f$ is provided and every function has a formula of size at most the length of its truth table (see Theorem 13).
When describing a function $f$ that is an input to MFSP, one naturally wants to denote by $n$ two different quantities: the number of variable inputs to a function $f$ and the length of the truth table of $f$ (which is the true input length for MFSP). We maintain the convention throughout this paper that $n$ denotes the input arity of $f$ and $N=2^n$ denotes the length of the truth table of $f$.
\subsection*{2.4 Useful Facts About Formulas}
We will make use of some basic facts about formulas in our work. First, one can easily bound the number of formulas of size at most $s$.
\begin{ntproposition}[Author Proposition 12]
The number of formulas on $n$-inputs of size at most $s$ is at most $2^{s\log n(1+o(1))}$.
\end{ntproposition}
We also know tight upper bounds on the maximum formula complexity of a $n$-input function.
\begin{nttheorem}[Author Theorem 13: Lozhkin {\cite{ref11}} improving on Lupanov {\cite{ref12}}]
Let $f:\{0,1\}^n\to\{0,1\}$. Then
\[L(f)\leq\frac{2^n}{\log n}\left(1+O\left(\frac 1{\log n}\right)\right).\]
\end{nttheorem}
Combining the size upper bound in Theorem 13 with the bound on the number of formulas of size $s$, we get the following proposition.
\begin{ntproposition}[Author Proposition 14: Random functions have not too many near optimal formulas]
Let $n$ and $k$ be positive integers. Let $N=2^n$. Assume $k=O(2^n/\log^2n)$. Then all but a $o(1)$-fraction of $n$-input Boolean functions $f$ satisfy
\[|\Canon{k}(f)|=2^{O(N/\log\log N)}.\]
\end{ntproposition}
\begin{proof}
Theorem 13 say that every $n$-input function has a formula of size at most $\frac{2^n}{\log n}(1+O(\frac 1{\log n}))$. Thus, any formula for computing $n$-input function that is within an additive $k$ of being optimal has size at most $s$ where
\[s\leq k+\frac{2^n}{\log n}(1+O(\tfrac 1{\log n}))=\frac{2^n}{\log n}(1+O(\tfrac 1{\log n})).\]
Proposition 12 implies that the number of formulas of size at most $s$ is upper bounded by
\[2^{s\log n(1+o(1))}=2^{N(1+O(1/\log\log N))}.\]
Hence, since there are $2^N$ Boolean functions on $n$-inputs, it follows that in expectation a random function has at most $2^{O(N/\log\log N)}$ formulas within $k$ of being optimal. The desired claim then follows by an application of Markov's inequality.
\end{proof}
We note that the bound given by Proposition 14 is actually counting formulas that are isomorphic to each other as distinct. Unfortunately removing this redundancy does not improve on the bound in Proposition 14. However, the fact that our results rely on the number of distinct formulas up to isomorphism means that there is no obvious obstruction to better bounds being proved and hence to our algorithms being more efficient.
We will also make use of the fact that integer comparison can be implemented by linear-sized formulas.
\begin{ntproposition}[Author Proposition 15: Small formulas for integer comparison]
Let $y\in\{0,1\}^n$. Let $\operatorname{GrtrThan}_y:\{0,1\}^n\to\{0,1\}$ be the function given by $\operatorname{GrtrThan}_y(x)=1$ if and only if $x>y$ in the usual lexicographic order on $\{0,1\}^n$. Then $L(\operatorname{GrtrThan}_y(x))\leq n$.
\end{ntproposition}
\begin{proof}
We work by induction on $n$. If $n=1$, then clearly $L(\operatorname{GrtrThan}_y)=1$ (either it is 0 if $y=1$ or it equals $x$ if $y=0$).
Now suppose $n>1$. Let $x_1,\ldots,x_n$ and $y_1,\ldots,y_n$ denote the bits of $x$ and $y$ respectively where $x_1$ and $y_1$ denotes the highest order bit. Let $x',y'\in\{0,1\}^{n-1}$ be given by $x'=x_2\ldots x_n$ and $y'=y_2\ldots y_n$ respectively.
Now, $x>y$ if and only if one of the following two statements is true:
\begin{itemize}\item $x_1>y_1$, or\item $x_1=y_1$ and $x'>y'$.\end{itemize}
Since $x_1,y_1\in\{0,1\}$, this is equivalent to
\[x>y\quad\Longleftrightarrow\quad(x_1>y_1)\lor(x'>y').\]
By induction this means $L(\operatorname{GrtrThan}_y)\leq 1+n-1=n$.
\end{proof}
\subsection*{2.5 Partial Functions and their Formula Size}
Partial functions will be a crucial building block in our reductions. A partial Boolean function is a function $\gamma:\{0,1\}^n\to\{0,1,?\}$ for some integer $n\geq 1$. We denote partial functions using Greek letters such as $\gamma$ and $\mu$, although sometimes we resort to the Roman alphabet with a$?$ subscript such as $h_?$.
In contrast, we say a Boolean function $f:\{0,1\}\to\{0,1\}$ is total Boolean function (though we allow for a partial Boolean function to indeed be total).
We say a total Boolean function $g$ agrees with a partial Boolean function $\gamma$ if
\[\gamma(x)\in\{0,1\}\quad\Longrightarrow\quad\gamma(x)=g(x).\]
One can naturally define the minimum formula size of a partial Boolean function $\gamma$ as follows
\[L(\gamma)=\min\{L(g):g\text{ is a total function that agrees with }\gamma\}.\]
The following theorem regarding the formula complexity of partial functions will be useful in our randomized worst-case reduction.
\begin{nttheorem}[Author Theorem 16: Pippenger {\cite{ref16}}]
Let $\gamma:\{0,1\}^n\to\{0,1,?\}$ be a partial function. Let $p_?=|\gamma^{-1}(?)|/2^n$. Then,
\[L(\gamma)\leq(1+o(1))\cdot(1-p_?)\frac{2^n}{\log n}.\]
\end{nttheorem}
